\documentclass[12pt,a4paper,twocolumn]{scrartcl}

\usepackage[english]{babel}
\usepackage[utf8x]{inputenc}
\usepackage{amsmath}
\usepackage{graphicx}
\usepackage{cite}
\usepackage{hyperref}
\usepackage{url}
\usepackage{float}
\usepackage{listings}
\usepackage[colorinlistoftodos]{todonotes}


\renewcommand{\baselinestretch}{1.5}

\begin{document}
\begin{titlepage}
    \centering
    \includegraphics[width=0.30\textwidth]{VUB.png}\par\vspace{1cm}
    {\scshape\Large Principles of Object-Oriented Programming Languages\par}
    \vspace{1cm}
    {\scshape\Large Assignement - Languages comparison : Python and Smalltalk\par}
    \vspace{1.5cm}
    {\Large\itshape Bruno Rocha Pereira - 0529512\par}
    \vfill
\end{titlepage}

\section{Closures}
\section{Delegation}
\section{Inheritance}
Python : Multiple inheritance
ST : Simple with methods to simulate multiple. No way to do an diamond inheritance
\section{Meta objects}
\section{Refactoring}
\section{Typing}
\section{Mixins}
Python has mixins.
Example : the socketserver module contains TCPServer and UDPServer as well as the ForkingMixin and ThreadingMixin mixins. Using a mixin with a class consists of creating a new class inheriting the class and the mixin. We can thus write a class ThreadingUDPServer:
class ThreadingUDPServer(ThreadingMixIn, UDPServer): pass
that will use threading to handle each connection to the UDP server.
Pharo has traits.

Defining a trait:
TAnnouncer>>announce: anAnnouncement
    | announcement |
    announcement := anAnnouncement asAnnouncement.
    self subscriptions ifNil: [ ^ announcement ].
    self subscriptions keysAndValuesDo: [ :class :actions|
        (class handles: announcement) 
            ifTrue: [ actions valueWithArguments: (Array with: announcement) ] ].
    ^ announcement
TAnnouncer>>on: anAnnouncementClass do: aValuable
    ^ self subscribe: anAnnouncementClass do: aValuable


Using it : 

Object subclass: #BulletinBoard
    uses: TAnnouncer
    instanceVariableNames: 'subscriptions'
    classVariableNames: ''
    poolDictionaries: ''
    category: 'BulletinBoard'

\footnote{http://pharo.gemtalksystems.com/book/LanguageAndLibraries/Traits/}
\section{Reflexion}
Inspection are allowed in both languages. 

>>> import os
>>> os.path = 42
>>> os.path
42
>>> os.path.join()
Traceback (most recent call last):
  File "<stdin>", line 1, in <module>
AttributeError: 'int' object has no attribute 'join'

However, in python some built-in functions (like print, datetime,..) cannot be changed (they can still be overwritten) since they are written in C.
The list of built-in functions can be found using dir(__builtins__)
Example of 
>>> print.__str__()
'<built-in function print>'
>>> print.__str__ = 42
Traceback (most recent call last):
  File "<stdin>", line 1, in <module>
AttributeError: 'builtin_function_or_method' object attribute '__str__' is read-only



\section{Concurrency}



\end{document}